% Dead-End Edge Case Diagram for Concave Hull Algorithm
% Import into main document with: % Dead-End Edge Case Diagram for Concave Hull Algorithm
% Import into main document with: % Dead-End Edge Case Diagram for Concave Hull Algorithm
% Import into main document with: % Dead-End Edge Case Diagram for Concave Hull Algorithm
% Import into main document with: \input{dead_end_diagram.tex}
% Requires: \usepackage{tikz}, \usetikzlibrary{calc}

\begin{tikzpicture}[
    scale=1.0,
    hull point/.style={circle, fill=red!70!black, inner sep=2pt},
    data point/.style={circle, fill=blue!40!black, inner sep=1.8pt},
    candidate point/.style={circle, fill=green!60!black, inner sep=2.2pt},
    current point/.style={circle, fill=orange!90!black, inner sep=3pt},
    hull edge/.style={thick, red!70!black},
    candidate edge/.style={dashed, green!50!black, thick},
    intersect marker/.style={red!80!black, thick},
]

% --- Hull points (the walk so far) ---
\coordinate (A) at (0, 0);       % firstPoint / start
\coordinate (B) at (1.4, 0.6);
\coordinate (C) at (2.8, 1.8);
\coordinate (D) at (4.0, 3.2);
\coordinate (E) at (3.2, 4.6);   % hull curves back inward
\coordinate (F) at (1.6, 4.2);
\coordinate (G) at (0.6, 3.0);   % hull narrows into a channel
\coordinate (H) at (1.8, 2.4);   % current point — deep inside

% --- Draw hull edges ---
\draw[hull edge] (A) -- (B) node[midway, below, font=\scriptsize, text=red!70!black] {$e_1$};
\draw[hull edge] (B) -- (C) node[midway, below right, font=\scriptsize, text=red!70!black] {$e_2$};
\draw[hull edge] (C) -- (D) node[midway, right, font=\scriptsize, text=red!70!black] {$e_3$};
\draw[hull edge] (D) -- (E) node[midway, right, font=\scriptsize, text=red!70!black] {$e_4$};
\draw[hull edge] (E) -- (F) node[midway, above, font=\scriptsize, text=red!70!black] {$e_5$};
\draw[hull edge] (F) -- (G) node[midway, left, font=\scriptsize, text=red!70!black] {$e_6$};
\draw[hull edge] (G) -- (H) node[midway, below left, font=\scriptsize, text=red!70!black] {$e_7$};

% --- Hull point markers ---
\node[hull point, label={[font=\scriptsize]below left:$p_1$}] at (A) {};
\node[hull point] at (B) {};
\node[hull point] at (C) {};
\node[hull point] at (D) {};
\node[hull point] at (E) {};
\node[hull point] at (F) {};
\node[hull point] at (G) {};
\node[current point, label={[font=\scriptsize, text=orange!90!black]below:$p_\text{cur}$}] at (H) {};

% --- Candidate points (k=3 nearest neighbours) ---
\coordinate (c1) at (-0.2, 1.6);
\coordinate (c2) at (0.4, 4.8);
\coordinate (c3) at (3.6, 1.2);

\node[candidate point, label={[font=\scriptsize]left:$c_1$}] at (c1) {};
\node[candidate point, label={[font=\scriptsize]above left:$c_2$}] at (c2) {};
\node[candidate point, label={[font=\scriptsize]right:$c_3$}] at (c3) {};

% --- Candidate edges (all intersect existing hull edges) ---
\draw[candidate edge] (H) -- (c1);
\draw[candidate edge] (H) -- (c2);
\draw[candidate edge] (H) -- (c3);

% --- Intersection markers (X) ---
% c1 candidate crosses e1 (approximately)
\coordinate (x1) at (0.65, 1.2);
\draw[intersect marker] ($(x1)+(-0.12,-0.12)$) -- ($(x1)+(0.12,0.12)$);
\draw[intersect marker] ($(x1)+(-0.12,0.12)$) -- ($(x1)+(0.12,-0.12)$);

% c2 candidate crosses e5 or e6 (approximately)
\coordinate (x2) at (1.1, 3.5);
\draw[intersect marker] ($(x2)+(-0.12,-0.12)$) -- ($(x2)+(0.12,0.12)$);
\draw[intersect marker] ($(x2)+(-0.12,0.12)$) -- ($(x2)+(0.12,-0.12)$);

% c3 candidate crosses e2 or e3 (approximately)
\coordinate (x3) at (2.95, 1.65);
\draw[intersect marker] ($(x3)+(-0.12,-0.12)$) -- ($(x3)+(0.12,0.12)$);
\draw[intersect marker] ($(x3)+(-0.12,0.12)$) -- ($(x3)+(0.12,-0.12)$);

% --- Scatter some background data points ---
\foreach \px/\py in {
    0.3/1.0, 0.8/2.2, 1.2/1.4, 2.0/0.5, 3.4/2.6,
    2.6/3.8, 1.0/4.6, 0.2/3.8, 3.8/4.0, 4.2/1.8,
    2.2/4.4, -0.4/2.4, 0.6/0.3, 3.0/0.4, 1.4/3.6,
    2.4/1.0, 3.6/3.4, 0.0/4.2
}{
    \node[data point] at (\px, \py) {};
}

% --- Legend ---
\begin{scope}[shift={(5.5, 4.5)}]
    \node[anchor=north west, font=\small\bfseries] at (0, 0.4) {Legend};
    \draw[hull edge] (0, 0) -- (0.5, 0);
    \node[anchor=west, font=\scriptsize] at (0.6, 0) {Hull edge};
    
    \draw[candidate edge] (0, -0.5) -- (0.5, -0.5);
    \node[anchor=west, font=\scriptsize] at (0.6, -0.5) {Candidate edge};
    
    \draw[intersect marker] (-0.1, -0.9) -- (0.1, -1.1);
    \draw[intersect marker] (-0.1, -1.1) -- (0.1, -0.9);
    \node[anchor=west, font=\scriptsize] at (0.6, -1.0) {Intersection};
    
    \node[current point] at (0.15, -1.5) {};
    \node[anchor=west, font=\scriptsize] at (0.6, -1.5) {Current point};
    
    \node[candidate point] at (0.15, -2.0) {};
    \node[anchor=west, font=\scriptsize] at (0.6, -2.0) {Candidate ($k$-NN)};
    
    \node[data point] at (0.15, -2.5) {};
    \node[anchor=west, font=\scriptsize] at (0.6, -2.5) {Data point};
\end{scope}

% --- Annotation ---
\node[anchor=north, font=\footnotesize, text width=5.5cm, align=center, text=black!70]
    at (2.2, -0.8) {%
    \textbf{Dead end:} The hull walk has curved inward,
    leaving $p_\text{cur}$ trapped. Every candidate edge
    from $p_\text{cur}$ to a $k$-NN intersects an existing
    hull edge, so no legal next step exists.%
};

\end{tikzpicture}
% Requires: \usepackage{tikz}, \usetikzlibrary{calc}

\begin{tikzpicture}[
    scale=1.0,
    hull point/.style={circle, fill=red!70!black, inner sep=2pt},
    data point/.style={circle, fill=blue!40!black, inner sep=1.8pt},
    candidate point/.style={circle, fill=green!60!black, inner sep=2.2pt},
    current point/.style={circle, fill=orange!90!black, inner sep=3pt},
    hull edge/.style={thick, red!70!black},
    candidate edge/.style={dashed, green!50!black, thick},
    intersect marker/.style={red!80!black, thick},
]

% --- Hull points (the walk so far) ---
\coordinate (A) at (0, 0);       % firstPoint / start
\coordinate (B) at (1.4, 0.6);
\coordinate (C) at (2.8, 1.8);
\coordinate (D) at (4.0, 3.2);
\coordinate (E) at (3.2, 4.6);   % hull curves back inward
\coordinate (F) at (1.6, 4.2);
\coordinate (G) at (0.6, 3.0);   % hull narrows into a channel
\coordinate (H) at (1.8, 2.4);   % current point — deep inside

% --- Draw hull edges ---
\draw[hull edge] (A) -- (B) node[midway, below, font=\scriptsize, text=red!70!black] {$e_1$};
\draw[hull edge] (B) -- (C) node[midway, below right, font=\scriptsize, text=red!70!black] {$e_2$};
\draw[hull edge] (C) -- (D) node[midway, right, font=\scriptsize, text=red!70!black] {$e_3$};
\draw[hull edge] (D) -- (E) node[midway, right, font=\scriptsize, text=red!70!black] {$e_4$};
\draw[hull edge] (E) -- (F) node[midway, above, font=\scriptsize, text=red!70!black] {$e_5$};
\draw[hull edge] (F) -- (G) node[midway, left, font=\scriptsize, text=red!70!black] {$e_6$};
\draw[hull edge] (G) -- (H) node[midway, below left, font=\scriptsize, text=red!70!black] {$e_7$};

% --- Hull point markers ---
\node[hull point, label={[font=\scriptsize]below left:$p_1$}] at (A) {};
\node[hull point] at (B) {};
\node[hull point] at (C) {};
\node[hull point] at (D) {};
\node[hull point] at (E) {};
\node[hull point] at (F) {};
\node[hull point] at (G) {};
\node[current point, label={[font=\scriptsize, text=orange!90!black]below:$p_\text{cur}$}] at (H) {};

% --- Candidate points (k=3 nearest neighbours) ---
\coordinate (c1) at (-0.2, 1.6);
\coordinate (c2) at (0.4, 4.8);
\coordinate (c3) at (3.6, 1.2);

\node[candidate point, label={[font=\scriptsize]left:$c_1$}] at (c1) {};
\node[candidate point, label={[font=\scriptsize]above left:$c_2$}] at (c2) {};
\node[candidate point, label={[font=\scriptsize]right:$c_3$}] at (c3) {};

% --- Candidate edges (all intersect existing hull edges) ---
\draw[candidate edge] (H) -- (c1);
\draw[candidate edge] (H) -- (c2);
\draw[candidate edge] (H) -- (c3);

% --- Intersection markers (X) ---
% c1 candidate crosses e1 (approximately)
\coordinate (x1) at (0.65, 1.2);
\draw[intersect marker] ($(x1)+(-0.12,-0.12)$) -- ($(x1)+(0.12,0.12)$);
\draw[intersect marker] ($(x1)+(-0.12,0.12)$) -- ($(x1)+(0.12,-0.12)$);

% c2 candidate crosses e5 or e6 (approximately)
\coordinate (x2) at (1.1, 3.5);
\draw[intersect marker] ($(x2)+(-0.12,-0.12)$) -- ($(x2)+(0.12,0.12)$);
\draw[intersect marker] ($(x2)+(-0.12,0.12)$) -- ($(x2)+(0.12,-0.12)$);

% c3 candidate crosses e2 or e3 (approximately)
\coordinate (x3) at (2.95, 1.65);
\draw[intersect marker] ($(x3)+(-0.12,-0.12)$) -- ($(x3)+(0.12,0.12)$);
\draw[intersect marker] ($(x3)+(-0.12,0.12)$) -- ($(x3)+(0.12,-0.12)$);

% --- Scatter some background data points ---
\foreach \px/\py in {
    0.3/1.0, 0.8/2.2, 1.2/1.4, 2.0/0.5, 3.4/2.6,
    2.6/3.8, 1.0/4.6, 0.2/3.8, 3.8/4.0, 4.2/1.8,
    2.2/4.4, -0.4/2.4, 0.6/0.3, 3.0/0.4, 1.4/3.6,
    2.4/1.0, 3.6/3.4, 0.0/4.2
}{
    \node[data point] at (\px, \py) {};
}

% --- Legend ---
\begin{scope}[shift={(5.5, 4.5)}]
    \node[anchor=north west, font=\small\bfseries] at (0, 0.4) {Legend};
    \draw[hull edge] (0, 0) -- (0.5, 0);
    \node[anchor=west, font=\scriptsize] at (0.6, 0) {Hull edge};
    
    \draw[candidate edge] (0, -0.5) -- (0.5, -0.5);
    \node[anchor=west, font=\scriptsize] at (0.6, -0.5) {Candidate edge};
    
    \draw[intersect marker] (-0.1, -0.9) -- (0.1, -1.1);
    \draw[intersect marker] (-0.1, -1.1) -- (0.1, -0.9);
    \node[anchor=west, font=\scriptsize] at (0.6, -1.0) {Intersection};
    
    \node[current point] at (0.15, -1.5) {};
    \node[anchor=west, font=\scriptsize] at (0.6, -1.5) {Current point};
    
    \node[candidate point] at (0.15, -2.0) {};
    \node[anchor=west, font=\scriptsize] at (0.6, -2.0) {Candidate ($k$-NN)};
    
    \node[data point] at (0.15, -2.5) {};
    \node[anchor=west, font=\scriptsize] at (0.6, -2.5) {Data point};
\end{scope}

% --- Annotation ---
\node[anchor=north, font=\footnotesize, text width=5.5cm, align=center, text=black!70]
    at (2.2, -0.8) {%
    \textbf{Dead end:} The hull walk has curved inward,
    leaving $p_\text{cur}$ trapped. Every candidate edge
    from $p_\text{cur}$ to a $k$-NN intersects an existing
    hull edge, so no legal next step exists.%
};

\end{tikzpicture}
% Requires: \usepackage{tikz}, \usetikzlibrary{calc}

\begin{tikzpicture}[
    scale=1.0,
    hull point/.style={circle, fill=red!70!black, inner sep=2pt},
    data point/.style={circle, fill=blue!40!black, inner sep=1.8pt},
    candidate point/.style={circle, fill=green!60!black, inner sep=2.2pt},
    current point/.style={circle, fill=orange!90!black, inner sep=3pt},
    hull edge/.style={thick, red!70!black},
    candidate edge/.style={dashed, green!50!black, thick},
    intersect marker/.style={red!80!black, thick},
]

% --- Hull points (the walk so far) ---
\coordinate (A) at (0, 0);       % firstPoint / start
\coordinate (B) at (1.4, 0.6);
\coordinate (C) at (2.8, 1.8);
\coordinate (D) at (4.0, 3.2);
\coordinate (E) at (3.2, 4.6);   % hull curves back inward
\coordinate (F) at (1.6, 4.2);
\coordinate (G) at (0.6, 3.0);   % hull narrows into a channel
\coordinate (H) at (1.8, 2.4);   % current point — deep inside

% --- Draw hull edges ---
\draw[hull edge] (A) -- (B) node[midway, below, font=\scriptsize, text=red!70!black] {$e_1$};
\draw[hull edge] (B) -- (C) node[midway, below right, font=\scriptsize, text=red!70!black] {$e_2$};
\draw[hull edge] (C) -- (D) node[midway, right, font=\scriptsize, text=red!70!black] {$e_3$};
\draw[hull edge] (D) -- (E) node[midway, right, font=\scriptsize, text=red!70!black] {$e_4$};
\draw[hull edge] (E) -- (F) node[midway, above, font=\scriptsize, text=red!70!black] {$e_5$};
\draw[hull edge] (F) -- (G) node[midway, left, font=\scriptsize, text=red!70!black] {$e_6$};
\draw[hull edge] (G) -- (H) node[midway, below left, font=\scriptsize, text=red!70!black] {$e_7$};

% --- Hull point markers ---
\node[hull point, label={[font=\scriptsize]below left:$p_1$}] at (A) {};
\node[hull point] at (B) {};
\node[hull point] at (C) {};
\node[hull point] at (D) {};
\node[hull point] at (E) {};
\node[hull point] at (F) {};
\node[hull point] at (G) {};
\node[current point, label={[font=\scriptsize, text=orange!90!black]below:$p_\text{cur}$}] at (H) {};

% --- Candidate points (k=3 nearest neighbours) ---
\coordinate (c1) at (-0.2, 1.6);
\coordinate (c2) at (0.4, 4.8);
\coordinate (c3) at (3.6, 1.2);

\node[candidate point, label={[font=\scriptsize]left:$c_1$}] at (c1) {};
\node[candidate point, label={[font=\scriptsize]above left:$c_2$}] at (c2) {};
\node[candidate point, label={[font=\scriptsize]right:$c_3$}] at (c3) {};

% --- Candidate edges (all intersect existing hull edges) ---
\draw[candidate edge] (H) -- (c1);
\draw[candidate edge] (H) -- (c2);
\draw[candidate edge] (H) -- (c3);

% --- Intersection markers (X) ---
% c1 candidate crosses e1 (approximately)
\coordinate (x1) at (0.65, 1.2);
\draw[intersect marker] ($(x1)+(-0.12,-0.12)$) -- ($(x1)+(0.12,0.12)$);
\draw[intersect marker] ($(x1)+(-0.12,0.12)$) -- ($(x1)+(0.12,-0.12)$);

% c2 candidate crosses e5 or e6 (approximately)
\coordinate (x2) at (1.1, 3.5);
\draw[intersect marker] ($(x2)+(-0.12,-0.12)$) -- ($(x2)+(0.12,0.12)$);
\draw[intersect marker] ($(x2)+(-0.12,0.12)$) -- ($(x2)+(0.12,-0.12)$);

% c3 candidate crosses e2 or e3 (approximately)
\coordinate (x3) at (2.95, 1.65);
\draw[intersect marker] ($(x3)+(-0.12,-0.12)$) -- ($(x3)+(0.12,0.12)$);
\draw[intersect marker] ($(x3)+(-0.12,0.12)$) -- ($(x3)+(0.12,-0.12)$);

% --- Scatter some background data points ---
\foreach \px/\py in {
    0.3/1.0, 0.8/2.2, 1.2/1.4, 2.0/0.5, 3.4/2.6,
    2.6/3.8, 1.0/4.6, 0.2/3.8, 3.8/4.0, 4.2/1.8,
    2.2/4.4, -0.4/2.4, 0.6/0.3, 3.0/0.4, 1.4/3.6,
    2.4/1.0, 3.6/3.4, 0.0/4.2
}{
    \node[data point] at (\px, \py) {};
}

% --- Legend ---
\begin{scope}[shift={(5.5, 4.5)}]
    \node[anchor=north west, font=\small\bfseries] at (0, 0.4) {Legend};
    \draw[hull edge] (0, 0) -- (0.5, 0);
    \node[anchor=west, font=\scriptsize] at (0.6, 0) {Hull edge};
    
    \draw[candidate edge] (0, -0.5) -- (0.5, -0.5);
    \node[anchor=west, font=\scriptsize] at (0.6, -0.5) {Candidate edge};
    
    \draw[intersect marker] (-0.1, -0.9) -- (0.1, -1.1);
    \draw[intersect marker] (-0.1, -1.1) -- (0.1, -0.9);
    \node[anchor=west, font=\scriptsize] at (0.6, -1.0) {Intersection};
    
    \node[current point] at (0.15, -1.5) {};
    \node[anchor=west, font=\scriptsize] at (0.6, -1.5) {Current point};
    
    \node[candidate point] at (0.15, -2.0) {};
    \node[anchor=west, font=\scriptsize] at (0.6, -2.0) {Candidate ($k$-NN)};
    
    \node[data point] at (0.15, -2.5) {};
    \node[anchor=west, font=\scriptsize] at (0.6, -2.5) {Data point};
\end{scope}

% --- Annotation ---
\node[anchor=north, font=\footnotesize, text width=5.5cm, align=center, text=black!70]
    at (2.2, -0.8) {%
    \textbf{Dead end:} The hull walk has curved inward,
    leaving $p_\text{cur}$ trapped. Every candidate edge
    from $p_\text{cur}$ to a $k$-NN intersects an existing
    hull edge, so no legal next step exists.%
};

\end{tikzpicture}
% Requires: \usepackage{tikz}, \usetikzlibrary{calc}

\begin{tikzpicture}[
    scale=1.0,
    hull point/.style={circle, fill=red!70!black, inner sep=2pt},
    data point/.style={circle, fill=blue!40!black, inner sep=1.8pt},
    candidate point/.style={circle, fill=green!60!black, inner sep=2.2pt},
    current point/.style={circle, fill=orange!90!black, inner sep=3pt},
    hull edge/.style={thick, red!70!black},
    candidate edge/.style={dashed, green!50!black, thick},
    intersect marker/.style={red!80!black, thick},
]

% --- Hull points (the walk so far) ---
\coordinate (A) at (0, 0);       % firstPoint / start
\coordinate (B) at (1.4, 0.6);
\coordinate (C) at (2.8, 1.8);
\coordinate (D) at (4.0, 3.2);
\coordinate (E) at (3.2, 4.6);   % hull curves back inward
\coordinate (F) at (1.6, 4.2);
\coordinate (G) at (0.6, 3.0);   % hull narrows into a channel
\coordinate (H) at (1.8, 2.4);   % current point — deep inside

% --- Draw hull edges ---
\draw[hull edge] (A) -- (B) node[midway, below, font=\scriptsize, text=red!70!black] {$e_1$};
\draw[hull edge] (B) -- (C) node[midway, below right, font=\scriptsize, text=red!70!black] {$e_2$};
\draw[hull edge] (C) -- (D) node[midway, right, font=\scriptsize, text=red!70!black] {$e_3$};
\draw[hull edge] (D) -- (E) node[midway, right, font=\scriptsize, text=red!70!black] {$e_4$};
\draw[hull edge] (E) -- (F) node[midway, above, font=\scriptsize, text=red!70!black] {$e_5$};
\draw[hull edge] (F) -- (G) node[midway, left, font=\scriptsize, text=red!70!black] {$e_6$};
\draw[hull edge] (G) -- (H) node[midway, below left, font=\scriptsize, text=red!70!black] {$e_7$};

% --- Hull point markers ---
\node[hull point, label={[font=\scriptsize]below left:$p_1$}] at (A) {};
\node[hull point] at (B) {};
\node[hull point] at (C) {};
\node[hull point] at (D) {};
\node[hull point] at (E) {};
\node[hull point] at (F) {};
\node[hull point] at (G) {};
\node[current point, label={[font=\scriptsize, text=orange!90!black]below:$p_\text{cur}$}] at (H) {};

% --- Candidate points (k=3 nearest neighbours) ---
\coordinate (c1) at (-0.2, 1.6);
\coordinate (c2) at (0.4, 4.8);
\coordinate (c3) at (3.6, 1.2);

\node[candidate point, label={[font=\scriptsize]left:$c_1$}] at (c1) {};
\node[candidate point, label={[font=\scriptsize]above left:$c_2$}] at (c2) {};
\node[candidate point, label={[font=\scriptsize]right:$c_3$}] at (c3) {};

% --- Candidate edges (all intersect existing hull edges) ---
\draw[candidate edge] (H) -- (c1);
\draw[candidate edge] (H) -- (c2);
\draw[candidate edge] (H) -- (c3);

% --- Intersection markers (X) ---
% c1 candidate crosses e1 (approximately)
\coordinate (x1) at (0.65, 1.2);
\draw[intersect marker] ($(x1)+(-0.12,-0.12)$) -- ($(x1)+(0.12,0.12)$);
\draw[intersect marker] ($(x1)+(-0.12,0.12)$) -- ($(x1)+(0.12,-0.12)$);

% c2 candidate crosses e5 or e6 (approximately)
\coordinate (x2) at (1.1, 3.5);
\draw[intersect marker] ($(x2)+(-0.12,-0.12)$) -- ($(x2)+(0.12,0.12)$);
\draw[intersect marker] ($(x2)+(-0.12,0.12)$) -- ($(x2)+(0.12,-0.12)$);

% c3 candidate crosses e2 or e3 (approximately)
\coordinate (x3) at (2.95, 1.65);
\draw[intersect marker] ($(x3)+(-0.12,-0.12)$) -- ($(x3)+(0.12,0.12)$);
\draw[intersect marker] ($(x3)+(-0.12,0.12)$) -- ($(x3)+(0.12,-0.12)$);

% --- Scatter some background data points ---
\foreach \px/\py in {
    0.3/1.0, 0.8/2.2, 1.2/1.4, 2.0/0.5, 3.4/2.6,
    2.6/3.8, 1.0/4.6, 0.2/3.8, 3.8/4.0, 4.2/1.8,
    2.2/4.4, -0.4/2.4, 0.6/0.3, 3.0/0.4, 1.4/3.6,
    2.4/1.0, 3.6/3.4, 0.0/4.2
}{
    \node[data point] at (\px, \py) {};
}

% --- Legend ---
\begin{scope}[shift={(5.5, 4.5)}]
    \node[anchor=north west, font=\small\bfseries] at (0, 0.4) {Legend};
    \draw[hull edge] (0, 0) -- (0.5, 0);
    \node[anchor=west, font=\scriptsize] at (0.6, 0) {Hull edge};
    
    \draw[candidate edge] (0, -0.5) -- (0.5, -0.5);
    \node[anchor=west, font=\scriptsize] at (0.6, -0.5) {Candidate edge};
    
    \draw[intersect marker] (-0.1, -0.9) -- (0.1, -1.1);
    \draw[intersect marker] (-0.1, -1.1) -- (0.1, -0.9);
    \node[anchor=west, font=\scriptsize] at (0.6, -1.0) {Intersection};
    
    \node[current point] at (0.15, -1.5) {};
    \node[anchor=west, font=\scriptsize] at (0.6, -1.5) {Current point};
    
    \node[candidate point] at (0.15, -2.0) {};
    \node[anchor=west, font=\scriptsize] at (0.6, -2.0) {Candidate ($k$-NN)};
    
    \node[data point] at (0.15, -2.5) {};
    \node[anchor=west, font=\scriptsize] at (0.6, -2.5) {Data point};
\end{scope}

% --- Annotation ---
\node[anchor=north, font=\footnotesize, text width=5.5cm, align=center, text=black!70]
    at (2.2, -0.8) {%
    \textbf{Dead end:} The hull walk has curved inward,
    leaving $p_\text{cur}$ trapped. Every candidate edge
    from $p_\text{cur}$ to a $k$-NN intersects an existing
    hull edge, so no legal next step exists.%
};

\end{tikzpicture}